\section{Low relations and simple factors}\label{sec:low-relations}

At stage \(i\) of the construction, we seek a sum of \(b\)-th powers
that lies in the degree-\(b\) ideal piece allowed by the new triangular
row, where \(b=a_i\). The earlier rows have degrees \(a_d\ge b\);
if we retained them only as equations of those degrees, they need not
contribute to this ideal piece. We instead retain them as consequences
of linear and quadratic relations among an enlarged collection of
formal linear forms.

We first describe the equations on a tuple of summands. Each is obtained
by summing the values of a degree-\(b\) polynomial \(w\) on those
summands. The next section requires every nonzero \(w\) in this linear
system to have an irreducible factor of multiplicity one.
We obtain this condition by excluding curves of bounded degree from
the projective zero set of the relations, and then construct initial
data for which this exclusion holds.

\subsection{The induction data}

Fix the regular sequence from Theorem~\ref{thm:construction}, and put
\[
 D=\bigl(\max_{1\le j\le n}a_j\bigr)^2.
\]
At the beginning of stage \(i\), where \(n\ge i\ge1\), the data are:
\begin{enumerate}[label=\textup{(\roman*)}]
\item a finitely generated field extension \(k/\C\) and a
finite-dimensional central division \(k\)-algebra \(\Delta_0\);
\item a finite-dimensional \(k\)-vector space \(U\), with
\(\dim U\ge3\), containing the formal variables \(x_1,\ldots,x_n\),
and a linear evaluation map
\[
 U\longrightarrow\Delta_0[x_1,\ldots,x_n]_1
\]
whose images commute pairwise and which sends each \(x_j\) to the
corresponding central variable;
\item a homogeneous ideal \(Q\subseteq\Sym_k U\), generated in
degrees one and two, which vanishes under evaluation, such that
\(\Sym_k U/Q\) is finite over \(k[x_1,\ldots,x_n]\);
\item formal elements \(t_d\in U\), for \(d>i\), whose rows
\eqref{eq:triangular} hold in \(\Sym_k U/Q\), with nonzero
coefficients \(c_d\in k\).
\end{enumerate}
In addition, after every algebraically closed field extension, the
closed set
\[
 \mathcal C(Q)=V_+(Q)\subseteq\PP(U^\vee)
\]
contains no integral projective curve of degree at most \(D\).
We call \(\mathcal C(Q)\) the characteristic set of \(Q\).
No smoothness, reducedness, or irreducibility is required of this set.
Polynomial degree in \(\Sym U\) always refers to the formal linear
forms, all of which have degree one.

We use the abstract complex embeddings of Section~\ref{subsec:conventions}
only for topology and numerical comparison; division is always proved
over the current field \(k\).

\subsection{The equations for the next row}

Assume the induction data at stage \(i\), and write \(b=a_i\).
The ideal relevant to the new power identity is
\begin{equation}\label{eq:stage-ideal}
 \mathfrak j=(Q,f_l\ (1\le l\le i),t_d\ (d>i))
       \subseteq\Sym_k U.
\end{equation}
Membership in \(\mathfrak j_b\subseteq\Sym^b U\) means that a
degree-\(b\) form has the shape required on the right side of the new
row, modulo \(Q\), with some scalar coefficient of \(f_i\).
That coefficient need not yet be nonzero. Section~\ref{sec:assembly}
will make it nonzero when choosing the sum of powers.

\begin{lemma}\label{lem:stage-basepoint}
The space \(\mathfrak j_b\) is base-point-free on \(\PP(U^\vee)\).
\end{lemma}

\begin{proof}
Every specified generator of \(\mathfrak j\) has degree at most
\(b\): the generators of \(Q\) have degree at most two,
\(b\ge2\), and \(a_l\le a_i=b\) for \(l\le i\).
At a projective point, a nonzero generator can be multiplied by a
monomial of complementary degree nonzero at that point.
Thus a base point of \(\mathfrak j_b\) would annihilate every
generator of \(\mathfrak j\).

At such a point \(f_l=0\) for \(l\le i\) and \(t_d=0\) for \(d>i\).
The earlier triangular rows, taken in the order
\(i+1,i+2,\ldots,n\), imply successively that the remaining \(f_d\)
also vanish. Equation~\eqref{eq:regular-hilbert} forces all \(x_j\)
to vanish, which is impossible at a point of \(\mathcal C(Q)\)
by finiteness over \(k[x]\).
\end{proof}

For \(\lambda\in(\Sym^b U)^\vee\), define the polynomial on \(U\)
by \(w_\lambda(\ell)=\lambda(\ell^b)\), and put
\begin{equation}\label{eq:stage-annihilator}
 W=\{\,w_\lambda:\lambda|_{\mathfrak j_b}=0\,\}.
\end{equation}
Characteristic zero identifies these functionals with degree-\(b\)
polynomials on \(U\), since pure powers span \(\Sym^b U\).
For any tuple \(\ell_1,\ldots,\ell_M\in U\), we have
\[
 \sum_{m=1}^M\ell_m^b\in\mathfrak j_b
 \quad\Longleftrightarrow\quad
 \sum_{m=1}^M w(\ell_m)=0
 \quad\text{for every }w\in W.
\]
These are the equations on the tuple of summands. To obtain an integral
parameter space with a highly connected smooth cone, Section~\ref{sec:parameters}
will use the condition that each nonzero \(w\in W\) has a factor of
multiplicity one. The next lemma derives that condition from the
absence of small curves in \(\mathcal C(Q)\).

\begin{lemma}\label{lem:simple-factor}
Let \(U\) be a vector space of dimension at least three over a
characteristic-zero field, and let \(b\ge2\).
Suppose that \(Q\subseteq\Sym U\) is generated in degrees one
and two, and that its geometric characteristic set contains no
integral curve of degree at most \(b(b-1)\).
Let \(\mathfrak j_b\subseteq\Sym^b U\) be a base-point-free
subspace containing \((Q)_b\), and define \(W\) by
\eqref{eq:stage-annihilator}.
Then every nonzero element of \(W\), after any algebraically closed
extension, has an irreducible factor of multiplicity one.
\end{lemma}

\begin{proof}
Fix \(w\in W\setminus\{0\}\) over an algebraically closed extension.
We first show that every repeated factor of \(w\) must be linear.
If there were no simple factor, a second gradient-map argument would
then force \(w\) to be a pure power, contradicting base-point-freeness.

For every linear or quadratic generator \(q\) of \(Q\),
the annihilation of
\(q\,\Sym^{b-\deg q}U\) gives the constant-coefficient
differential equation
\begin{equation}\label{eq:annihilator-differential}
 q(\partial)w=0.
\end{equation}
Here a basis of \(U\) identifies its elements with the corresponding
constant differential operators on polynomial functions on \(U\).
This is the differential-operator description of the degreewise
annihilator pairing; see \cite[pp.~1082--1083, Definitions~1--2
and~(5a)]{EmsalemIarrobino95}.

Suppose an irreducible factor \(p\) occurs in \(w=p^r h\) with
\(r\ge2\) and \(p\nmid h\). For a quadratic \(q\), the term of
lowest possible \(p\)-adic order in \(q(\partial)w\) is
\[
 r(r-1)p^{r-2}h\,q(\nabla p)
       \pmod{p^{r-1}}.
\]
For a linear \(q\), it is
\(rp^{r-1}h\,q(\nabla p)\pmod{p^r}\).
Equation~\eqref{eq:annihilator-differential} therefore shows that,
generically on the hypersurface \(p=0\), its projective gradient
belongs to \(\mathcal C(Q)\).

If that gradient map is nonconstant, it has nonzero differential
at some smooth point \(a\) of \(p=0\). Choose a tangent vector \(v\)
at \(a\) detected by this differential, and an ambient tangent vector
\(v_{\perp}\) transverse to the hypersurface. A projective plane through \(a\)
with tangent plane spanned by \(v,v_{\perp}\) cuts the hypersurface in a curve
smooth at \(a\), with tangent line spanned by \(v\). Its component
through \(a\) has degree at most \(\deg p\), and its gradient map is
nonconstant. On the normalization of this component, the partial
derivatives are sections of the pullback of \(\cO(\deg p-1)\).
After removing their common base divisor, they define a morphism
whose pulled-back hyperplane bundle has degree at most
\(\deg p(\deg p-1)\). This degree is the degree of the image curve
times the degree of the morphism onto its image. The integral image
is therefore a curve in \(\mathcal C(Q)\) of degree at most
\(b(b-1)\), a contradiction.
If the gradient map is constant instead, Euler's identity
places the hypersurface \(p=0\) in the hyperplane perpendicular
to that constant gradient. Irreducibility then forces \(p\) to
be linear. Thus every repeated irreducible factor of \(w\) is linear.

Assume, seeking a contradiction, that \(w\) has no simple factor.
We may then write
\[
 w=c\prod_{\nu=1}^r p_\nu^{m_\nu},
 \qquad m_\nu\ge2,
\]
with distinct linear forms \(p_\nu\), and put
\(v_\nu=\nabla p_\nu\in U^\vee\).
The preceding argument gives \([v_\nu]\in\mathcal C(Q)\).
For a quadratic generator \(q\), logarithmic differentiation yields
\[
 \frac{q(\partial)w}{w}
 =
 q\left(\sum_{\nu=1}^r\frac{m_\nu v_\nu}{p_\nu}\right)
 -\sum_{\nu=1}^r\frac{m_\nu q(v_\nu)}{p_\nu^2}.
\]
The left side and each \(q(v_\nu)\) vanish. Together with the
linear differential equations, this proves that the ambient
projective gradient map of \(w\) has image in \(\mathcal C(Q)\).
Here the map is defined on the ambient \(\PP(U)\), not just on
one factor hypersurface. If it is nonconstant, restrict it to a line
along which it varies. The partial derivatives restrict to sections
of \(\cO_{\PP^1}(b-1)\); removing their common zeros gives a
nonconstant morphism whose image has degree at most \(b-1\).
That image is again a forbidden curve in \(\mathcal C(Q)\).

Finally, a constant projective gradient means that all partial
derivatives of \(w\) are proportional to one fixed vector.
After a linear change of coordinates, \(w\) depends on only one
coordinate. Homogeneity gives \(w=c\,p^b\) for a nonzero linear
form \(p\in U^\vee\). The corresponding functional on \(\Sym^b U\)
is, up to the nonzero scalar \(c\), evaluation at \(p\).
Its annihilation of \(\mathfrak j_b\) contradicts base-point freeness.
\end{proof}

The geometric hypothesis has now supplied the required factor property
of \(W\). It remains to produce initial division-algebra data with this
hypothesis and to show that the finite enlargements in the induction
preserve it.

\subsection{Starting with generic square roots}

We construct the initial relations as \(s_\rho^2=q_\rho(x)\), with
independent generic quadrics \(q_\rho\). A curve in their projective
zero set would give a finite cover of a curve in \(\PP^{n-1}\), on
which all the \(q_\rho\) acquire square roots. For a curve of bounded
degree, sufficiently many of these square roots remain independent,
forcing the cover degree to exceed the bound. The following field-of-definition
estimate makes this argument uniform over all such curves.

\begin{lemma}\label{lem:bounded-curve-field}
For integers \(n\ge2\) and \(D\ge1\), there is an integer \(d_*(n,D)\)
with the following property. If \(\Omega\) is an algebraically closed
extension of \(\C\) and
\(Z\subseteq\PP^{n-1}_\Omega\) is an integral curve of degree at most
\(D\), then \(Z\) is defined over an algebraically closed intermediate
field \(F_0\subseteq\Omega\) with
\(\operatorname{trdeg}_\C F_0\le d_*(n,D)\).
\end{lemma}

\begin{proof}
We use the hyperplane-incidence description of the Chow form
\cite[Section~1, Proposition~1.1]{Guerra96}, recalling the argument for
curves to make the field-of-definition bound explicit.
Consider pairs of hyperplanes whose intersection meets \(Z\).
The incidence variety over \(Z\), with the two hyperplanes chosen
through a point of \(Z\), is irreducible of dimension \(2n-3\).
A general pair in its image meets \(Z\) in finitely many points,
so its image is an irreducible divisor in
\((\PP^{n-1})^\vee\times(\PP^{n-1})^\vee\).
Its defining bihomogeneous equation has degree \(\deg Z\) in each
block: for a general first hyperplane, the corresponding divisor
in the second block is the reduced union of the \(\deg Z\)
hyperplanes representing the points of intersection.

This equation recovers \(Z\). A point on \(Z\) has the property that
every pair of hyperplanes through it satisfies the equation.
For a point outside \(Z\), choose a first hyperplane through it
meeting \(Z\) properly, and then a second hyperplane through it
avoiding the resulting finite set. That pair does not satisfy the
equation. This recovery condition is defined over the field of coefficients.
Indeed, on a point chart \(x_a\ne0\), a hyperplane through \(x\) has
coefficient \(h_a=-\sum_{j\ne a}h_jx_j/x_a\). Substitute this expression
and its counterpart for the second hyperplane into the incidence
equation, clear denominators, and set every coefficient in the remaining
hyperplane variables to zero. These finite equations express the recovery
condition on the chart and agree on overlaps. Over the algebraic closure
of the coefficient field, take the reduced structure. In characteristic
zero it is geometrically reduced, so extension to \(\Omega\) gives
exactly \(Z\).

An equation of bidegree \((d,d)\), \(d\le D\), uses at most
\(\binom{n-1+D}{D}^2\) coefficients. Adjoin those coefficients to
\(\C\), then take its algebraic closure inside \(\Omega\).
This proves the claim, for instance with
\(d_*(n,D)=\binom{n-1+D}{D}^2\).
\end{proof}

\begin{proposition}\label{prop:initial-data}
For every \(n\ge1\) and \(D\ge1\), there exist data
\(k,\Delta_0,U,Q\) satisfying conditions \textup{(i)}--\textup{(iii)}
above and the condition on \(\mathcal C(Q)\).
They may be chosen with
\[
 U=\langle x_1,\ldots,x_n,s_1,\ldots,s_m\rangle_k,
 \qquad
 Q=(s_\rho^2-q_\rho(x):1\le\rho\le m),
\]
where \(m\ge2\) and the coefficient vectors of the quadrics \(q_\rho\)
are jointly algebraically independent over \(\C\).
\end{proposition}

\begin{proof}
Choose an even integer \(v_0\) and ordinary linear forms
\(L_1,\ldots,L_{v_0}\in\C[x]_1\) whose squares span \(\C[x]_2\).
Such a choice is possible by polarization, increasing \(v_0\) if
necessary. For an integer \(m\) to be specified, form the skew Laurent
algebra on symbols
\[
 y_{\rho j}^{\pm1},
 \qquad 1\le\rho\le m,\quad 1\le j\le v_0,
\]
in which distinct generators in the same row anticommute and
generators in different rows commute. Ordered Laurent monomials
give a basis: multiplication is defined by reordering and inserting
the prescribed signs. Their leading terms show that this algebra
is a domain.

The squares of the generators are central and algebraically
independent. Localize their Laurent polynomial algebra to
\[
 k=\C(y_{\rho j}^2:1\le\rho\le m,\ 1\le j\le v_0).
\]
The resulting algebra \(\Delta_0\) is a finite-dimensional domain
over \(k\), with basis the monomials whose exponents are zero or one;
it is therefore a division algebra. Its center is exactly \(k\).
Indeed, a parity vector \(\alpha\) in one row commutes with its
\(j\)-th generator precisely when
\(\sum_{l\ne j}\alpha_l=0\) modulo two. Writing
\(T=\sum_l\alpha_l\), these equations give \(\alpha_j=T\) for all
\(j\). Since \(v_0\) is even, they imply \(T=0\) and then
\(\alpha=0\). Applying this in every row proves the claim; linear
independence of the parity monomials prevents cancellations.

Evaluate the new formal forms by
\[
 s_\rho=\sum_{j=1}^{v_0}y_{\rho j}L_j(x).
\]
Forms from different rows commute, and their squares are
\[
 s_\rho^2=q_\rho(x)
 :=\sum_{j=1}^{v_0}y_{\rho j}^2L_j(x)^2.
\]
The full-rank linear map from the independent square variables
to the coefficients of each quadric shows that the coefficient
vectors of the \(q_\rho\) are jointly algebraically independent.
The indicated ideal \(Q\) consists of identities, and its quotient
is finite over \(k[x]\), since every \(s_\rho\) is integral.

This supplies the commuting forms and the quadratic relations.
We now choose \(m\) so that the resulting finite cover contains no
curve of degree at most \(D\).
Over any algebraically closed extension \(\Omega\) of \(k\), the
linear projection
\[
 \mathcal C(Q)_\Omega\longrightarrow\PP^{n-1}_\Omega,
 \qquad [x:s]\longmapsto[x],
\]
is defined everywhere and is finite. Its pullback of \(\cO(1)\)
is the ambient \(\cO(1)\). If \(n=1\), its target is a point,
and there is no curve upstairs. We may then take \(m=2\).

Assume \(n\ge2\), and suppose an integral curve \(T\) of degree at
most \(D\) occurs upstairs. Its image is an integral curve \(Z\);
if the map degree is \(d\), then
\[
 \deg T=d\deg Z.
\]
In particular \(d\le D\) and \(\deg Z\le D\).
Choose a field of definition \(F_0\subseteq\Omega\) for \(Z\)
as in Lemma~\ref{lem:bounded-curve-field}, and write \(d_*=d_*(n,D)\).
Let \(r=\binom{n+1}{2}\), the size of each coefficient block.
Since all \(mr\) coefficients are independent over \(\C\),
their transcendence degree over \(F_0\) is at least \(mr-d_*\).

At least \(m-d_*\) entire coefficient blocks are jointly independent
over \(F_0\). For completeness, scan the blocks, retaining a block
whenever it contributes all \(r\) new independent parameters over
the retained blocks and \(F_0\). Each rejected block contributes at
most \(r-1\) even after all retained blocks have been adjoined.
Indeed, enlarging the field generated by the retained blocks can only
decrease the transcendence degree contributed by a rejected block.
Adjoin all retained blocks first, then the rejected ones.
If \(j\) blocks were rejected, the total transcendence degree would
be at most \(mr-j\), so \(j\le d_*\).

Restrict the retained generic quadrics to \(Z\).
Their zero divisors are nonempty, reduced, supported in the smooth
locus, and pairwise disjoint. They can also be taken disjoint from
the zeros of a fixed coordinate \(x_j\) nonzero on \(Z\).
These are nonempty open conditions on the jointly generic full
quadrics: the complete linear system of quadrics is base-point-free,
separates tangent directions on the smooth curve, and can avoid
each specified finite set. At a zero of the \(\rho\)-th such divisor,
\(q_\rho/x_j^2\) has valuation one, while the other retained
functions have valuation zero. Thus their classes are independent
in \(\Omega(Z)^\times/\Omega(Z)^{\times2}\).
The function field of \(T\) contains square roots of all these
functions, by the equations \(s_\rho^2=q_\rho\).
Consequently
\[
 d=[\Omega(T):\Omega(Z)]\ge2^{m-d_*}
\]
by the elementary degree formula for independent squareclasses.
Choose \(m\ge2\) with \(2^{m-d_*}>D\). This contradicts \(d\le D\)
and proves the proposition.
\end{proof}

\subsection{Preserving the geometric condition}

The initial data now satisfy the required curve exclusion. Each later
stage will enlarge the formal relation algebra by a finite graded
extension, still generated in degree one. The following degree
comparison shows why the same exclusion survives.

\begin{lemma}\label{lem:sparsity-preserved}
Let \(U\subseteq\widetilde U\) be finite-dimensional vector spaces
over a field, and let \(Q\subseteq\Sym U\) and
\(\widetilde Q\subseteq\Sym\widetilde U\) be homogeneous ideals,
with \(Q\subseteq\widetilde Q\).
Suppose that \(\Sym\widetilde U/\widetilde Q\) is finite over
\(\Sym U/Q\). Then the induced projection of characteristic sets
is everywhere defined and finite, and pulls back \(\cO(1)\) to
\(\cO(1)\). In particular, if the first characteristic set contains
no integral curve of degree at most \(D\) after algebraically
closed extension, neither does the second.
\end{lemma}

\begin{proof}
The quotient of the larger graded algebra by all positive-degree
old forms is finite dimensional. Hence no point of its projective
scheme has all old coordinates zero. The finite graded algebra
map therefore induces the stated finite projective morphism,
with the asserted hyperplane bundle. On an integral curve, its
degree equals the degree of its image multiplied by the positive
degree of the restricted morphism. A forbidden curve upstairs
would thus give a forbidden curve downstairs.
\end{proof}

These arguments include \(b=2\), when a repeated factor of a quadratic
would already make it a pure square. They also include \(n=1\):
the auxiliary forms in Proposition~\ref{prop:initial-data} ensure
\(\dim U\ge3\), as required by the factor lemma.

We have constructed the initial induction data and shown that their
curve exclusion survives every finite homogeneous enlargement.
At each stage, \(D\ge b(b-1)\), so Lemmas~\ref{lem:stage-basepoint}
and~\ref{lem:simple-factor} show that every nonzero member of the
space \(W\) in \eqref{eq:stage-annihilator} has a simple factor.
Its equations \(\sum_m w(\ell_m)=0\) are exactly the conditions that
\(\sum_m\ell_m^b\) belong to \(\mathfrak j_b\).
The next section studies their parameter complete intersection.
